\section{Elliptic operators in exponential coordinates}\label{sec:g443}

The archimedean argument uses two local facts about an elliptic differential
operator. Analytic coefficients force its smooth null solutions to be analytic.
For smooth coefficients and order \(m>n/2\) in positive dimension \(n\),
a compactly supported kernel gives the value of a null solution at a point.
The kernel must be chosen before the solution.

\subsection{Coordinates and principal symbols}
Let \(n,m\) be nonnegative integers, let \(\Omega\subset\mathbb R^n\) be open, and let \(A\) be a
finite set of multi-indices of length \(|\alpha|\leq m\). The coefficients are
complex-valued functions of real variables. Write
\begin{equation}\label{eq:g443-operator}
 Pu(x)=\sum_{\alpha\in A}a_\alpha(x)\partial^\alpha u(x),\qquad
 p_m(x,\xi)=\sum_{\substack{\alpha\in A\\|\alpha|=m}}a_\alpha(x)\xi^\alpha .
\end{equation}
Ellipticity means \(p_m(x,\xi)\ne0\) for \(x\in\Omega\) and real
\(\xi\ne0\). With the convention \(\partial_j\mapsto i\xi_j\), the top
symbol is \(i^m p_m\), with the same characteristic set.

Here is the passage from the ordered derivatives used in a Lie chart to
\eqref{eq:g443-operator}. For a word \(J=(j_1,\ldots,j_k)\) in the coordinate
directions \(e_1,\ldots,e_n\), put
\[
 u_J=D^ku[e_{j_1},\ldots,e_{j_k}],\qquad
 \nu(J)_i=\#\{r:j_r=i\}.
\]
The empty word gives \(u\). For smooth \(u\) and arbitrary fixed vectors
\(v_1,\ldots,v_k\), the function \(D^ku[v_1,\ldots,v_k]\) is smooth.
Its derivative along \(v_0\) is \(D^{k+1}u[v_0,v_1,\ldots,v_k]\).
For a multi-index \(\alpha\), the coordinate list consists of
\(\alpha_i\) copies of \(e_i\), in increasing order of \(i\). Its length
is \(|\alpha|=\sum_i\alpha_i\). Evaluation on this list defines
\(\partial^\alpha u\), and \(\xi^\alpha=\prod_i\xi_i^{\alpha_i}\).
Symmetry permits any permutation, so \(\partial_i u_J=u_{(i,J)}\) and
\(u_J=\partial^{\nu(J)}u\).
For an ordered operator \(\sum_{k=0}^m\sum_{|J|=k}c_J u_J\), take
\[
 b_\alpha(x)=\sum_{\nu(J)=\alpha}c_J(x) \quad (|\alpha|\leq m).
\]
This finite regrouping preserves the operator on smooth functions and
preserves its degree-\(m\) symbol at every \(x,\xi\), since
\(\prod_{r=1}^m\xi_{j_r}=\xi^{\nu(J)}\). It preserves analyticity of the
coefficients on a neighbourhood of each point where all the \(c_J\) have
this property, even if the designated set of points is arbitrary. The
coefficients for different orderings of a word may be different.
For a vector field \(V=\sum_i v_i\partial_i\) the product rule gives
\[
 V(c_Ju_J)=(Vc_J)u_J+c_J\sum_i v_i u_{(i,J)} .
\]
Thus composition on the left by \(V\) multiplies the top symbol by
\(\sum_i v_i(x)\xi_i\). Derivatives of coefficients enter the lower orders.
This identity uses only the value of \(V\) at the point and the formula
\(L(v)=\sum_i v_i L(e_i)\) for a real-linear differential \(L\).
If the \(v_i\) and \(c_J\) are
analytic on a neighbourhood, the displayed coefficients are analytic there.

\subsection{The two local results}
\begin{proposition}[Interior analyticity]\label{prop:g443-analytic}
Suppose \(a_\alpha\) is real analytic on \(\Omega\) for every \(\alpha\in A\),
and \(P\) is elliptic there. If \(u\) is complex-valued, smooth on \(\Omega\)
and \(Pu=0\) on \(\Omega\), then \(u\) is real analytic on \(\Omega\).
\end{proposition}
\begin{proof}
If \(n=0\), every function on the domain is locally constant and hence analytic.
If \(n>0\) and \(m=0\),
ellipticity makes the multiplication coefficient nonzero at each point,
so \(u=0\) there. For \(m>0\), regard \(u\) as a distribution locally.
Analytic microlocal elliptic regularity gives
\[
 \mathrm{WF}_{\!A}(u)\subset
 \mathrm{WF}_{\!A}(Pu)\cup
 \{(x,\xi):x\in\Omega,\ \xi\ne0,\ p_m(x,\xi)=0\}
\]
by H\"ormander~\cite[Thm.~8.6.1]{g443-hormander}.
Both sets on the right are empty. By Theorem~8.4.5 of the same volume,
the projection of the analytic wavefront set is the analytic singular support.
The resulting analytic representative
equals the given smooth function pointwise.
\end{proof}

\begin{proposition}[Localized evaluation]\label{prop:g443-kernel}
Suppose \(n>0,\ m\geq1,\ n<2m\). Let the \(a_\alpha\) be smooth on
\(\Omega\), and \(P\) elliptic there. For every open \(U\subseteq\Omega\)
and \(x_0\in U\), there are complex-valued functions \(E,\beta\) on
\(\mathbb R^n\) with
\[
 E\in L^1(\mathbb R^n)\cap C^\infty(\mathbb R^n\setminus\{x_0\}),\quad
 \beta\in C^\infty_c(\mathbb R^n),\quad
 \operatorname{supp}E\ \cup\ \operatorname{supp}\beta\Subset U .
\]
For every \(\phi:\mathbb R^n\to\mathbb C\) whose restriction to \(U\) is
smooth,
\begin{equation}\label{eq:g443-eval}
 \int_{\mathbb R^n} E(x)P\phi(x)\,dx
       =\phi(x_0)-\int_{\mathbb R^n}\beta(x)\phi(x)\,dx .
\end{equation}
Both integrands are extended by zero outside \(U\).
\end{proposition}
\begin{proof}
Here the support of a function is the closure of its nonzero set, and
\(\Subset\) denotes compact containment.
The distribution pairing is complex bilinear. On a ball
\(x_0\in V\Subset U\) the transpose is
\[
 T={}^tP,\qquad Tv=\sum_{\alpha\in A}(-\partial)^\alpha(a_\alpha v).
\]
In the convention of \eqref{eq:g443-operator}, its principal symbol is
\((-1)^mp_m\). Continuity and ellipticity give a positive lower bound for
\(|p_m|\) on \(\overline V\) times the unit sphere. The properly supported parametrix theorem
on open sets gives \(Q\) of order \(-m\) with \(TQ=I-R\), where \(R\)
has a smooth kernel~\cite[Cor.~7.19]{g443-grubb}. Set
\(v=Q\delta_{x_0}\) and \(r=R\delta_{x_0}\). Then \(Tv=\delta_{x_0}-r\).
In \(V\setminus\{x_0\}\), smooth elliptic regularity makes \(v\) smooth.
Choose \(n/2<s<m\), possible by the dimension assumption. The Fourier
transform of a point mass has constant absolute value, so the integral of
\((1+|\xi|^2)^{-s}\) shows that \(\delta_{x_0}\in H^{-s}\). The local
Sobolev mapping rule gives \(v\in H_{\mathrm{loc}}^{m-s}(V)\subset
L^2_{\mathrm{loc}}(V)\)~\cite[\S7.4, proof of Cor.~7.20]{g443-grubb}.
Use its smooth representative off \(x_0\) and assign any finite value at
\(x_0\), a null set.

Take \(\eta\in C^\infty_c(V)\) equal to one near \(x_0\) and extend
\[
        E=\eta v,\qquad \beta=\eta r-[T,\eta]v
\]
by zero. Compact support and Cauchy--Schwarz give \(E\in L^1\).
The commutator is supported where \(\eta\) varies, away from \(x_0\),
so \(\beta\) is smooth with compact support in \(V\).
The identity \(TE=\eta Tv+[T,\eta]v=\delta_{x_0}-\beta\) gives
\eqref{eq:g443-eval} for compactly supported smooth test functions.
For the stated \(\phi\), take a cutoff in \(C^\infty_c(U)\) equal to one
near \(\{x_0\}\cup\operatorname{supp}E\cup\operatorname{supp}\beta\).
Its product with \(\phi\), extended by zero, is such a test function, and
all three terms of \eqref{eq:g443-eval} agree with those for \(\phi\).
\end{proof}

\subsection{Return to the convolution kernels}
In the application a fixed Lie algebra basis gives exponential coordinates
of dimension \(d\) on a neighbourhood of the identity. The chart
calculation supplies analytic coefficients and an elliptic principal symbol
for a monic polynomial in the Laplacian. The chart and its Haar density
are inputs supplied by the adjoining group-theoretic argument.
Proposition~\ref{prop:g443-analytic} gives analyticity at the origin of the
chart for each translated solution.

For a monic polynomial \(B\) of positive degree and \(d>0\), use
\(B^{d+1}\) in the evaluation formula. Its differential order is
\(m=2(d+1)\deg B\), so \(d<2m\). Every \(f\) killed by \(B(\Delta)\)
is killed by \(B(\Delta)^{d+1}\). The chart operator, its kernel \(\beta\),
and the chart itself can all be fixed before \(f\) or the translating element
\(g\) is chosen. Equation~\eqref{eq:g443-eval} for a translated solution
now gives \(u(0)=\int\beta u\).

Let \(b:\mathbb R^d\to\mathfrak g\) be the coordinate isomorphism and put
\(\chi_\pm(x)=\exp(\pm bx)\). Choose a ball about zero inside both ellipticity
neighbourhoods so that both charts map it into the prescribed identity neighbourhood.
Using the given Laplacian in each chart yields kernels \(\beta_\pm\).
Write \(\chi_+^*d\mu=j(x)\,dx\), where \(j\) is smooth and positive.
Define both group kernels in the positive chart by
\(\kappa_\pm(\chi_+(x))=\beta_\pm(x)/j(x)\) and extend them by zero.
Their supports are compactly contained in the chart, so the extensions are
smooth. Since \(\chi_-(x)=\chi_+(x)^{-1}\), for every such \(f\) and every \(g\)
\[
 f(g)=\int_G f(gy)\kappa_+(y)\,d\mu(y)
     =\int_G f(gy^{-1})\kappa_-(y)\,d\mu(y).
\]
In dimension zero the group is discrete and
\(\boldsymbol{1}_{\{1\}}/\mu(\{1\})\) gives both kernels.
